\chapter{Vector Analysis \& Differential Field Operators}
\index{Vector Analysis}
\index{Dot Product}
\index{Cross Product}
\index{Scalar Triple Product}
\index{Vector Triple Product}
\index{Gradient}
\index{Divergence}
\index{Curl}
\index{Laplacian}
\index{Conservative Fields}
\index{Lorentz Force}

\begin{chapterobjectives}
Upon mastering this chapter, students will be able to:
\begin{enumerate}
    \item Perform sophisticated algebraic manipulations with dot, cross, and triple vector products.
    \item Compute and physically interpret gradient ($\nabla$), divergence ($\nabla \cdot$), and curl ($\nabla \times$) in Cartesian and curvilinear geometries.
    \item Evaluate conservative vector fields, construct scalar potential functions, and test for irrotational and solenoidal conditions.
    \item Formulate vector identities essential for electrodynamics, continuum mechanics, and fluid dynamics.
    \item Visualize vector fields and streamline topologies computationally using scientific Python.
\end{enumerate}
\end{chapterobjectives}

\section{Vector Algebra and Triple Products}
Physical quantities characterized by both magnitude and spatial orientation are represented mathematically as vectors. In three-dimensional Euclidean space, let $\vec{A} = A_x\hat{i} + A_y\hat{j} + A_z\hat{k}$ and $\vec{B} = B_x\hat{i} + B_y\hat{j} + B_z\hat{k}$.

The scalar (dot) product is defined by:
\begin{equation}
\vec{A} \cdot \vec{B} = |\vec{A}||\vec{B}|\cos\theta = A_x B_x + A_y B_y + A_z B_z
\end{equation}
The vector (cross) product produces an orthogonal axial vector according to the right-hand rule:
\begin{equation}
\vec{A} \times \vec{B} = |\vec{A}||\vec{B}|\sin\theta\,\hat{n} = \begin{vmatrix} \hat{i} & \hat{j} & \hat{k} \\ A_x & A_y & A_z \\ B_x & B_y & B_z \end{vmatrix}
\end{equation}

\subsection{Triple Products and Geometric Interpretations}
The \textbf{scalar triple product} represents the signed volume of the parallelepiped subtended by vectors $\vec{A}, \vec{B}, \vec{C}$:
\begin{equation}
\vec{A} \cdot (\vec{B} \times \vec{C}) = \begin{vmatrix} A_x & A_y & A_z \\ B_x & B_y & B_z \\ C_x & C_y & C_z \end{vmatrix}
\end{equation}
This product is cyclic: $\vec{A} \cdot (\vec{B} \times \vec{C}) = \vec{B} \cdot (\vec{C} \times \vec{A}) = \vec{C} \cdot (\vec{A} \times \vec{B})$.

The \textbf{vector triple product} satisfies the celebrated ``BAC--CAB'' identity:
\begin{equation}
\vec{A} \times (\vec{B} \times \vec{C}) = \vec{B}(\vec{A} \cdot \vec{C}) - \vec{C}(\vec{A} \cdot \vec{B})
\end{equation}
This identity is fundamental in classical mechanics for formulating the angular momentum of rotating rigid bodies $\vec{L} = \sum m_i [\vec{r}_i \times (\vec{\omega} \times \vec{r}_i)]$.

\section{Differential Field Operators: Gradient, Divergence, and Curl}
Spatial variation of physical fields requires the Del (Nabla) vector differential operator:
\begin{equation}
\nabla \equiv \hat{i}\frac{\partial}{\partial x} + \hat{j}\frac{\partial}{\partial y} + \hat{k}\frac{\partial}{\partial z}
\end{equation}

\begin{mathdefinition}[The Gradient of a Scalar Field]
For a scalar field $\psi(x, y, z)$, the \textbf{gradient} $\nabla \psi$ is a vector field pointing in the direction of the greatest rate of spatial increase of $\psi$, with magnitude equal to that maximum directional rate of change:
\begin{equation}
\nabla \psi = \frac{\partial \psi}{\partial x}\hat{i} + \frac{\partial \psi}{\partial y}\hat{j} + \frac{\partial \psi}{\partial z}\hat{k}
\end{equation}
Geometrically, $\nabla \psi$ is perpendicular to the equipotential or level surface $\psi(x, y, z) = C$ at every point.
\end{mathdefinition}

\begin{mathdefinition}[The Divergence of a Vector Field]
The \textbf{divergence} of a vector field $\vec{A}$ represents the net volumetric flux density outflow per unit infinitesimal volume from a given point:
\begin{equation}
\nabla \cdot \vec{A} = \frac{\partial A_x}{\partial x} + \frac{\partial A_y}{\partial y} + \frac{\partial A_z}{\partial z}
\end{equation}
If $\nabla \cdot \vec{A} = 0$ everywhere, the field is called \textbf{solenoidal} (e.g., the magnetic field $\vec{B}$ in Maxwell's second law $\nabla \cdot \vec{B} = 0$, indicating the non-existence of magnetic monopoles).
\end{mathdefinition}

\begin{mathdefinition}[The Curl of a Vector Field]
The \textbf{curl} of a vector field $\vec{A}$ measures the microscopic circulation density or rotational vorticity around an infinitesimal loop:
\begin{equation}
\nabla \times \vec{A} = \left(\frac{\partial A_z}{\partial y} - \frac{\partial A_y}{\partial z}\right)\hat{i} + \left(\frac{\partial A_x}{\partial z} - \frac{\partial A_z}{\partial x}\right)\hat{j} + \left(\frac{\partial A_y}{\partial x} - \frac{\partial A_x}{\partial y}\right)\hat{k}
\end{equation}
A vector field with $\nabla \times \vec{A} = \vec{0}$ is called \textbf{irrotational} or \textbf{conservative}.
\end{mathdefinition}

\begin{maththeorem}[Null Identities of Vector Calculus]
For any twice continuously differentiable scalar field $\psi$ and vector field $\vec{A}$:
\begin{align}
\nabla \times (\nabla \psi) &= \vec{0} \quad \text{(The curl of any gradient field is identically zero)} \\
\nabla \cdot (\nabla \times \vec{A}) &= 0 \quad \text{(The divergence of any curl field is identically zero)}
\end{align}
\end{maththeorem}

\begin{pitfallbox}[Common Operator Confusion]
Do not confuse the Laplacian of a scalar with the vector Laplacian. 
For a scalar field:
\begin{equation}
\nabla^2 \psi = \nabla \cdot (\nabla \psi) = \frac{\partial^2 \psi}{\partial x^2} + \frac{\partial^2 \psi}{\partial y^2} + \frac{\partial^2 \psi}{\partial z^2}
\end{equation}
For a vector field, the vector Laplacian is defined via the curl-of-curl identity:
\begin{equation}
\nabla^2 \vec{A} \equiv \nabla(\nabla \cdot \vec{A}) - \nabla \times (\nabla \times \vec{A})
\end{equation}
In Cartesian coordinates, $\nabla^2 \vec{A} = (\nabla^2 A_x)\hat{i} + (\nabla^2 A_y)\hat{j} + (\nabla^2 A_z)\hat{k}$, but this component-wise equality fails in curvilinear coordinates because unit vectors depend on spatial coordinates!
\end{pitfallbox}

\section{Conservative Vector Fields and Potential Theory}
A vector force field $\vec{F}$ is conservative if the work done $W = \int_{C} \vec{F} \cdot d\vec{r}$ between any two points is independent of the path. The following four statements are mathematically equivalent in a simply connected domain:
\begin{enumerate}
    \item $\oint_C \vec{F} \cdot d\vec{r} = 0$ for every closed contour $C$.
    \item $\int_A^B \vec{F} \cdot d\vec{r}$ is path-independent.
    \item $\vec{F}$ can be expressed as the negative gradient of a scalar potential: $\vec{F} = -\nabla V$.
    \item $\vec{F}$ is irrotational everywhere: $\nabla \times \vec{F} = \vec{0}$.
\end{enumerate}

\section{Worked Examples and Physical Applications}

\begin{psctexample}[Testing Conservatism and Finding the Potential Function]
Consider the gravitational/electrostatic type force field:
\begin{equation}
\vec{F}(x, y, z) = (2xy + z^3)\hat{i} + x^2\hat{j} + 3xz^2\hat{k}
\end{equation}
Determine whether $\vec{F}$ is conservative. If so, construct its potential energy function $V(x, y, z)$ such that $\vec{F} = -\nabla V$ with $V(0,0,0) = 0$.

\textbf{\color{rbruNavy}Picture / Conceptualization:}
The field components depend on polynomial combinations of spatial coordinates. We test whether the field possesses microscopic vorticity.

\textbf{\color{rbruNavy}Strategy / Governing Laws:}
Compute $\nabla \times \vec{F}$. If the curl vanishes identically, determine $V$ by integrating the partial derivative conditions $-\frac{\partial V}{\partial x} = F_x$, $-\frac{\partial V}{\partial y} = F_y$, $-\frac{\partial V}{\partial z} = F_z$.

\textbf{\color{rbruNavy}Calculation / Analytical Solution:}
1. Evaluate the curl:
\begin{equation}
\nabla \times \vec{F} = \begin{vmatrix} \hat{i} & \hat{j} & \hat{k} \\ \frac{\partial}{\partial x} & \frac{\partial}{\partial y} & \frac{\partial}{\partial z} \\ 2xy + z^3 & x^2 & 3xz^2 \end{vmatrix}
= \hat{i}(0 - 0) - \hat{j}(3z^2 - 3z^2) + \hat{k}(2x - 2x) = \vec{0}
\end{equation}
Because $\nabla \times \vec{F} = \vec{0}$ everywhere, $\vec{F}$ is conservative.

2. Integrate for $V(x, y, z)$:
\begin{align}
\frac{\partial V}{\partial x} &= -(2xy + z^3) \implies V(x, y, z) = -x^2 y - xz^3 + g(y, z) \\
\frac{\partial V}{\partial y} &= -x^2 + \frac{\partial g}{\partial y} = -x^2 \implies \frac{\partial g}{\partial y} = 0 \implies g(y, z) = h(z) \\
\frac{\partial V}{\partial z} &= -3xz^2 + h'(z) = -3xz^2 \implies h'(z) = 0 \implies h(z) = C
\end{align}
Given the reference condition $V(0,0,0) = 0$, the integration constant is $C = 0$.
Thus:
\begin{equation}
V(x, y, z) = -(x^2 y + xz^3)
\end{equation}

\textbf{\color{rbruNavy}Think / Physical Insights:}
Because $\vec{F} = -\nabla V$, moving a test particle along any closed trajectory in this field yields net zero energy dissipation, ensuring complete conservation of mechanical energy $E = T + V$.
\qedexam
\end{psctexample}

\begin{pythonbox}[Vector Field Divergence and Streamline Visualization]
import numpy as np
import matplotlib.pyplot as plt

# Define 2D grid
Y, X = np.mgrid[-2:2:100j, -2:2:100j]

# Dipole-like electric vector field
U = X / (X**2 + Y**2 + 0.1)**1.5
V = Y / (X**2 + Y**2 + 0.1)**1.5

# Numerical divergence: dU/dx + dV/dy
dx = X[0, 1] - X[0, 0]
dy = Y[1, 0] - Y[0, 0]
div = np.gradient(U, axis=1)/dx + np.gradient(V, axis=0)/dy

print(f"Max divergence near core: {np.max(div):.2f}")
print("Streamlines smoothly diverge radially outwards from the positive source.")
\end{pythonbox}

\begin{psctexample}[Magnetic Torque on a Current Loop]
A rectangular current loop of width $a$ and height $b$ carries a steady current $I$. The loop lies in the $xy$-plane with its normal $\hat{n} = \hat{k}$. An external uniform magnetic field $\vec{B} = B_0(\hat{i} + \sqrt{3}\,\hat{j})$ is applied. Find (i) the net magnetic force on the loop, (ii) the magnetic torque $\vec{\tau} = \vec{m} \times \vec{B}$, and (iii) the potential energy $U = -\vec{m}\cdot\vec{B}$.

\textbf{Solution:}

The magnetic dipole moment is $\vec{m} = I(ab)\hat{k} = IAb\hat{k}$ where $A = ab$.

\textbf{(i) Net force on a rigid current loop in a uniform field is zero:}
\begin{equation}
\vec{F}_\text{net} = I\oint (d\vec{l} \times \vec{B}) = I\left(\oint d\vec{l}\right) \times \vec{B} = \vec{0}
\end{equation}
since $\oint d\vec{l} = \vec{0}$ for a closed path.

\textbf{(ii) Magnetic torque:}
\begin{align}
\vec{\tau} &= \vec{m} \times \vec{B} = (IAb\hat{k}) \times B_0(\hat{i} + \sqrt{3}\,\hat{j}) \\
           &= IAb B_0\left[(\hat{k}\times\hat{i}) + \sqrt{3}(\hat{k}\times\hat{j})\right] \\
           &= IAb B_0\left[\hat{j} + \sqrt{3}(-\hat{i})\right] = IAb B_0(-\sqrt{3}\,\hat{i} + \hat{j})
\end{align}
The torque magnitude is:
\begin{equation}
\boxed{|\vec{\tau}| = IAb B_0\sqrt{3 + 1} = 2IAb B_0}
\end{equation}

\textbf{(iii) Potential energy:}
\begin{equation}
U = -\vec{m}\cdot\vec{B} = -(IAb\hat{k})\cdot B_0(\hat{i} + \sqrt{3}\,\hat{j}) = 0
\end{equation}
The potential energy is zero because $\vec{m} \perp \vec{B}$ (the loop is in a metastable orientation).

\textbf{\color{rbruNavy}Think / Physical Insights:}
The torque $\vec{\tau}$ acts to align $\vec{m}$ with $\vec{B}$, reducing $U$ to its minimum $-|\vec{m}||\vec{B}|$. This is the fundamental operating principle of electric motors and galvanometers. The torque is maximum when $\vec{m} \perp \vec{B}$ (as here) and vanishes at equilibrium.
\qedexam
\end{psctexample}

\section{Summary of Key Formulas}
\begin{table}[H]
\centering
\small
\begin{tabularx}{\textwidth}{llX}
\toprule
\textbf{Operator / Law} & \textbf{Mathematical Definition} & \textbf{Physical Meaning} \\
\midrule
Gradient $\nabla \psi$ & $\sum_i \frac{1}{h_i}\frac{\partial \psi}{\partial u_i}\hat{e}_i$ & Maximum directional spatial derivative; normal to level surfaces \\
Divergence $\nabla \cdot \vec{A}$ & $\frac{1}{h_1 h_2 h_3}\sum_i \frac{\partial}{\partial u_i}\left(\frac{h_1 h_2 h_3}{h_i}A_i\right)$ & Net outflow flux density per unit volume (sources and sinks) \\
Curl $\nabla \times \vec{A}$ & $\frac{1}{h_1 h_2 h_3}\begin{vmatrix} h_1\hat{e}_1 & h_2\hat{e}_2 & h_3\hat{e}_3 \\ \partial/\partial u_1 & \partial/\partial u_2 & \partial/\partial u_3 \\ h_1 A_1 & h_2 A_2 & h_3 A_3 \end{vmatrix}$ & Circulation vorticity around microscopic loops \\
Laplacian $\nabla^2 \psi$ & $\nabla \cdot (\nabla \psi)$ & Mean curvature; appears in Poisson and wave equations \\
Magnetic Torque & $\vec{\tau} = \vec{m} \times \vec{B}$, $\;U = -\vec{m}\cdot\vec{B}$ & Aligns current loop with external field (motor principle) \\
\bottomrule
\end{tabularx}
\caption{Summary of differential vector operators in orthogonal curvilinear coordinates.}
\label{tab:ch02_summary}
\end{table}

\section{Chapter Exercises}
\begin{enumerate}
    \item \textbf{BAC-CAB Expansion:} Prove algebraically that $(\vec{A} \times \vec{B}) \cdot (\vec{C} \times \vec{D}) = (\vec{A} \cdot \vec{C})(\vec{B} \cdot \vec{D}) - (\vec{A} \cdot \vec{D})(\vec{B} \cdot \vec{C})$.
    \item \textbf{Angular Velocity Field:} The velocity field of a rigid body rotating with constant angular velocity $\vec{\omega} = \omega\hat{k}$ is $\vec{v} = \vec{\omega} \times \vec{r}$. Calculate $\nabla \cdot \vec{v}$ and $\nabla \times \vec{v}$. What is the physical significance of $\nabla \times \vec{v}$?
    \item \textbf{Solenoidal and Irrotational Check:} Show that the field $\vec{B}(\vec{r}) = \frac{\vec{m} \times \vec{r}}{r^3}$ (magnetic dipole field) satisfies $\nabla \cdot \vec{B} = 0$ for $r \neq 0$.
    \item \textbf{Orbital Angular Momentum:} An electron in a hydrogen atom moves in a circular orbit of radius $a_0 = 0.529\;\text{\AA}$ with speed $v = 2.19 \times 10^6\;\text{m/s}$. Calculate the orbital angular momentum $\vec{L} = m_e\vec{r} \times \vec{v}$ and show that $|\vec{L}| = \hbar$ (the reduced Planck constant).
    \item \textbf{Stokes' Theorem Application:} The vector field $\vec{F} = (y^2 - z^2)\hat{i} + (z^2 - x^2)\hat{j} + (x^2 - y^2)\hat{k}$. Use Stokes' theorem to evaluate $\oint_C \vec{F} \cdot d\vec{l}$ where $C$ is the intersection of the sphere $x^2 + y^2 + z^2 = 1$ with the plane $x + y + z = 0$.
\end{enumerate}
